\section{Metodología}
\label{sec:metodologia}

\subsection{Números aleatorios basados en contador}

Un programa paralelo puede ser rápido y estar equivocado. En Monte Carlo, un flujo aleatorio sembrado por hilo o por proceso haría depender el resultado del reparto de trabajo y del orden de planificación, y un error de ese tipo produce cifras de apariencia normal. Por eso se adoptó Philox4x32-10~\cite{salmon2011}, un generador basado en contador: cada bloque de cuatro palabras de 32 bits es una función pura $\text{philox}(\text{contador}, \text{clave})$ de diez rondas, sin estado interno. La implementación propia coincide con los vectores de prueba de Random123.

\begin{itemize}
  \item \textbf{Clave} (64 bits): se deriva de la semilla global $s$ y del identificador de configuración $c$ con dos pasadas de SplitMix64; no depende de hilos ni de procesos.
  \item \textbf{Contador}: $\{i_{\text{bajo}}, i_{\text{alto}}, d, \text{sub}\}$, donde $i = c\cdot N + j$ es el índice global del escenario, $d$ el día y $\text{sub}$ el uso: $\text{sub} = 0$ decide la ocurrencia y $\text{sub} \ge 1$ los intentos del muestreo gamma.
\end{itemize}

Cada número aleatorio queda determinado por $(s, c, i, d, \text{sub})$. Junto con tres hechos de la implementación, esto garantiza resultados idénticos bit a bit en todos los backends: (1) cada escenario consume los mismos números, lo procese quien lo procese; (2) los cálculos en punto flotante ($A_t$, $\Amax$) ocurren dentro de un escenario y siempre en el mismo orden, y ningún valor real se reduce entre hilos ni procesos; y (3) lo que sí se reduce ($k$, $n$ y las 512 clases del histograma) son enteros sin signo de 64 bits, cuya suma es asociativa y conmutativa.

\subsection{Protocolo de verificación}

La Tabla~\ref{tab:validaciones} lista las pruebas V1--V8. V1, V3, V4 y V7 forman parte de la suite CTest que la integración continua (GitHub Actions) ejecuta en cada cambio, con 1, 2, 4, 8 y 16 procesos en el caso MPI. V5 se ejecuta, también en la integración continua, con ThreadSanitizer~\cite{serebryany2009} y Archer~\cite{atzeni2016}, el detector de carreras de LLVM para OpenMP, con un control negativo que introduce una carrera deliberada y debe ser detectada.

\begin{table*}[!t]
\centering
\caption{Pruebas de verificación y estado al cierre del borrador}
\label{tab:validaciones}
\begin{tabular}{@{}llp{7.4cm}l@{}}
\toprule
Id & Prueba & Criterio de aceptación & Estado \\
\midrule
V1 & Equivalencia CPU & $k$, $n$ e histograma idénticos bit a bit entre secuencial, OpenMP, MPI E1/E2 e híbrido, para cualquier número de hilos y procesos & Aprobada \\
V2 & Equivalencia GPU & Diferencia relativa del conteo $\le 10^{-6}$ frente a la CPU & Pendiente (CUDA en desarrollo) \\
V3 & Conservación & $\sum$ escenarios procesados $= N$ y $\sum$ histograma $= n$ & Aprobada \\
V4 & Reproducibilidad & Dos corridas con la misma configuración dan resultados idénticos & Aprobada \\
V5 & Ausencia de carreras & Cero reportes de ThreadSanitizer + Archer; el control negativo sí se detecta & Aprobada \\
V6 & Validación estadística & Error relativo $\le 10$\,\% en frecuencia de lluvia, media, desviación, P95, P99, racha media y máximo diario; prueba K--S no rechazada al 5\,\% & Parcial (ver texto) \\
V7 & Caso analítico & Con $D=1$, $P = p_{01}\,Q\!\left(\alpha, (C+U)/\theta\right)$ dentro del IC de Wilson & Aprobada \\
V8 & Memoria y carreras GPU & Cero errores de \codigo{compute-sanitizer} & Pendiente (CUDA en desarrollo) \\
\bottomrule
\end{tabular}
\end{table*}

En V7, con un solo día se tiene $\Amax = \max(0, R - C)$, por lo que $\Amax > U \iff R > C + U$; la probabilidad exacta es la de llover el primer día ($p_{01}$, con el día previo seco) por la cola superior de la gamma, donde $Q$ es la función gamma incompleta regularizada superior.

V6 tiene dos partes. Frente a sus propios parámetros, el generador reproduce los siete estadísticos con errores relativos inferiores al \num{0.2}\,\% y la prueba de Kolmogorov--Smirnov no se rechaza ($N = 10^{6}$ temporadas). Frente a las series observadas del SENAMHI, en cambio, la frecuencia de días lluviosos, la media de los montos y la racha media se reproducen dentro del 10\,\%, pero los estadísticos de cola lo superan: la desviación de los montos falla en los 16 grupos estación--régimen auditados, el P99 en 15, el P95 en 12 y el máximo diario por temporada en 11. Por ejemplo, en Lambayeque (El Niño) el P99 simulado es \num{44.4}\,\% menor que el observado. La gamma de un solo parámetro de forma, ajustada por Thom, subestima la cola de la lluvia observada. Esto implica que las probabilidades de inundación serán, previsiblemente, una cota inferior en las capacidades de drenaje altas, donde dominan los eventos extremos.
% TODO-V6: la Ola D investiga el fallo de V6 observado (docs/validacion-v6.md). Sustituir este párrafo por el diagnóstico final: causa, corrección aplicada (si la hay) y tabla de errores por estadístico.

\subsection{Protocolo de medición}

Todos los puntos comparables usan la misma carga $N$, los mismos parámetros y la misma semilla. Cada punto se mide con una corrida de calentamiento, que no entra en la media, y cinco repeticiones; se informan la media y la desviación estándar muestral. Antes de calcular un speedup se comprueba V1 sobre cada corrida, calentamiento incluido: si $k$, $n$ o el histograma difieren de la referencia secuencial, la medición se descarta. Se rechazan las compilaciones \codigo{Debug}. Cada fila del CSV registra commit, CPU, núcleos físicos y lógicos, memoria, compilador, tipo de compilación, afinidad, hilos, procesos, distribución por nodo y fecha.

En OpenMP se fijan \codigo{OMP\_PROC\_BIND=close} y \codigo{OMP\_PLACES=cores} y se retiran \codigo{OMP\_NUM\_THREADS} y \codigo{OMP\_SCHEDULE} del entorno, de modo que hilos y \emph{schedule} solo provienen de la línea de comandos. Se comparan 2, 4, 8 y 16 recursos CPU por backend: hilos en OpenMP, procesos en MPI y núcleos totales $P\times T$ en el híbrido. Un punto cuyo número de hilos o núcleos totales supera los núcleos físicos del equipo se marca como sobresuscrito y se analiza aparte. El tiempo $T_p$ es el de pared del cómputo y la reducción; en MPI se mide con \codigo{MPI\_Wtime} en el rango 0 entre una barrera inicial y el final de la reducción, y el arranque de contenedores se excluye.

\subsection{Métricas}

Con $T_1$ el tiempo secuencial y $T_p$ el tiempo con $p$ recursos:
\begin{align}
S_p &= \frac{T_1}{T_p}, \qquad E_p = \frac{S_p}{p}, \qquad \Theta_p = \frac{N}{T_p}, \label{eq:speedup}\\
e_p &= \frac{1/S_p - 1/p}{1 - 1/p}, \label{eq:karpflatt}
\end{align}
donde $\Theta_p$ es el throughput (escenarios por segundo) y $e_p$ la fracción serial determinada experimentalmente de Karp y Flatt~\cite{karp1990}. Si $e_p$ es aproximadamente constante al crecer $p$, la pérdida de eficiencia se explica por una fracción serial fija, como en la ley de Amdahl~\cite{amdahl1967}; si crece con $p$, indica sobrecarga creciente (comunicación, sincronización o desequilibrio). En escalabilidad débil, donde la carga crece con los recursos, la referencia es la ley de Gustafson~\cite{gustafson1988}. La desviación del speedup se aproxima propagando los errores relativos de $T_1$ y $T_p$; no es un intervalo de confianza.

Se definen tres experimentos de rendimiento: X1, escalabilidad fuerte con carga fija; X2, escalabilidad débil con carga proporcional a los recursos; y X3, barrido de carga $N = 10^{6}, \dots, 10^{9}$ para estimar el peso del costo fijo. X4 compara los esquemas MPI estático y dinámico con los mismos procesos en nodos heterogéneos.
